% !TeX spellcheck = en_US
\documentclass[12pt,a4paper]{scrartcl}

\usepackage[utf8]{inputenc}
\usepackage[T1]{fontenc}
\usepackage{mathpazo}  % Palatino
\usepackage{geometry}
\geometry{top=2.3cm, bottom=1.3cm, left=1.5cm, right=1.5cm}

\usepackage{amsmath, amssymb}
\usepackage{graphicx, caption, subcaption}
\usepackage{booktabs}
\usepackage{siunitx}
\usepackage{listings}
\usepackage{hyperref}
\hypersetup{
	colorlinks=true,
	linkcolor=black,
	urlcolor=blue,
	citecolor=black
}

\usepackage{fancyhdr}
\usepackage{xcolor}
\usepackage{textcomp}
\usepackage{ragged2e}
\usepackage{algorithm}
\usepackage{algpseudocode}

% --- Colores personalizados ---
\definecolor{matlabgreen}{rgb}{0.0,0.6,0.0}
\definecolor{matlabblue}{rgb}{0.0,0.0,0.8}
\definecolor{lightgray}{gray}{0.97}
\definecolor{gray}{gray}{0.4}
\definecolor{orange}{rgb}{1,0.5,0}
\definecolor{purple}{rgb}{0.58,0,0.82}

% --- Estilo para MATLAB ---
\lstdefinestyle{matlabstyle}{
	language=Matlab,
	backgroundcolor=\color{lightgray!10},
	basicstyle=\ttfamily\footnotesize,
	keywordstyle=\color{blue}\bfseries,
	stringstyle=\color{purple},
	commentstyle=\color{matlabgreen}\itshape,
	numberstyle=\tiny\color{gray},
	numbers=left,
	stepnumber=1,
	showstringspaces=false,
	breaklines=true,
	breakatwhitespace=true,
	tabsize=4,
	frame=single,
	rulecolor=\color{gray!50},
	morekeywords={function,end,if,else,elseif,for,while,switch,case,otherwise,return,break,continue,try,catch},
}

\pagestyle{fancy}
\fancyhead[L]{Technical Report - Path Planning}
\fancyhead[R]{\thepage}
\fancyfoot{}

\title{
    \begin{center}
        \textbf{Path Planning with RRT* Algorithm \\ for Skid-Steer Mobile Robot} \\
        \vspace{0.3cm}
        {\large Diploma in Robotics - Planning and Perception Module}
    \end{center}
}

\author{Brayan Gerson Duran Toconas}
\date{December, 2024}

\begin{document}
	
	\maketitle
	\tableofcontents
	\newpage
	
	% ============================================================================
	\section{Introduction}
	% ============================================================================
	
	\justify
	This report presents the implementation of a path planning system using the \textbf{RRT* (Rapidly-exploring Random Tree Star)} algorithm for autonomous navigation of a skid-steer mobile robot. The system generates collision-free waypoints that are followed by an existing Fuzzy+PID+Lag Compensator controller.
	
	The implementation was developed in MATLAB and integrates seamlessly with the robot control system developed in the Control Laboratory module.
	
	% ============================================================================
	\section{Path Planning Method: RRT*}
	% ============================================================================
	
	\subsection{Algorithm Selection}
	
	The \textbf{RRT* (Rapidly-exploring Random Tree Star)} algorithm was selected for this implementation. RRT* is an extension of the basic RRT algorithm that provides \textit{asymptotic optimality}, meaning it converges to the optimal path as the number of iterations increases.
	
	\subsection{Reasons for Selection}
	
	The following reasons justify the selection of RRT* over other path planning algorithms:
	
	\begin{enumerate}
		\item \textbf{Asymptotic Optimality}: Unlike basic RRT, which finds \textit{any} feasible path, RRT* continuously optimizes the path through its rewiring mechanism, converging to the optimal solution.
		
		\item \textbf{Complex Environment Handling}: RRT* efficiently handles environments with multiple obstacles of various shapes (circles, rectangles, mazes) without requiring a grid discretization.
		
		\item \textbf{Rewiring Mechanism}: The algorithm continuously improves the tree structure by rewiring nodes when a shorter path is found, resulting in smoother and more efficient paths.
		
		\item \textbf{Probabilistic Completeness}: If a path exists, RRT* will find it given sufficient iterations.
		
		\item \textbf{Anytime Algorithm}: The algorithm can be stopped at any time and still provide a valid (though possibly suboptimal) path.
	\end{enumerate}
	
	\subsection{Comparison with Other Algorithms}
	
	\begin{table}[htb]
		\centering
		\caption{Comparison of Path Planning Algorithms}
		\begin{tabular}{lccc}
			\toprule
			\textbf{Algorithm} & \textbf{Optimality} & \textbf{Continuous Space} & \textbf{Complex Obstacles} \\
			\midrule
			Dijkstra / A* & Optimal (grid) & No & Limited \\
			Potential Fields & Not guaranteed & Yes & Poor (local minima) \\
			RRT (basic) & Not optimal & Yes & Good \\
			\textbf{RRT*} & \textbf{Asymptotically Optimal} & \textbf{Yes} & \textbf{Excellent} \\
			\bottomrule
		\end{tabular}
		\label{tab:comparison}
	\end{table}
	
	% ============================================================================
	\section{Algorithm Description}
	% ============================================================================
	
	\subsection{RRT* Pseudocode}
	
	\begin{algorithm}[H]
		\caption{RRT* Algorithm}
		\begin{algorithmic}[1]
			\State Initialize tree $T$ with start node $x_{start}$
			\For{$i = 1$ to $N_{max}$}
				\State $x_{rand} \gets$ SampleRandom() \Comment{With goal bias}
				\State $x_{nearest} \gets$ NearestNode($T$, $x_{rand}$)
				\State $x_{new} \gets$ Steer($x_{nearest}$, $x_{rand}$, $\eta$)
				\If{CollisionFree($x_{nearest}$, $x_{new}$)}
					\State $X_{near} \gets$ NearNodes($T$, $x_{new}$, $r$)
					\State $x_{min} \gets$ ChooseParent($X_{near}$, $x_{new}$) \Comment{Minimum cost}
					\State AddNode($T$, $x_{new}$, $x_{min}$)
					\State Rewire($T$, $X_{near}$, $x_{new}$) \Comment{Optimize tree}
					\If{Distance($x_{new}$, $x_{goal}$) $<$ $\epsilon$}
						\State MarkGoalReached()
					\EndIf
				\EndIf
			\EndFor
			\State \Return ExtractPath($T$, $x_{goal}$)
		\end{algorithmic}
	\end{algorithm}
	
	\subsection{Key Components}
	
	\subsubsection{Random Sampling with Goal Bias}
	The algorithm samples random points in the configuration space. A \textit{goal bias} parameter (default 10\%) causes the algorithm to occasionally sample the goal position directly, accelerating convergence.
	
	\subsubsection{Nearest Neighbor and Steering}
	For each random sample, the nearest node in the tree is found using Euclidean distance. The tree is extended from this node toward the sample by a maximum step size $\eta$.
	
	\subsubsection{Collision Checking}
	Before adding a new node, the algorithm verifies that the path segment is collision-free. This includes checking against:
	\begin{itemize}
		\item Circular obstacles (analytical distance check)
		\item Rectangular obstacles (discretized segment check)
		\item Robot radius safety margin
	\end{itemize}
	
	\subsubsection{Parent Selection and Rewiring}
	The key improvement of RRT* over RRT is the parent selection and rewiring:
	\begin{enumerate}
		\item Find all nodes within radius $r$ of the new node
		\item Choose the parent that minimizes the total cost from start
		\item Rewire neighboring nodes if connecting through the new node reduces their cost
	\end{enumerate}
	
	\subsection{Path Smoothing}
	
	After finding a path, a smoothing algorithm is applied:
	\begin{enumerate}
		\item \textbf{Shortcut Smoothing}: Randomly select pairs of non-adjacent waypoints and connect them directly if collision-free
		\item \textbf{Waypoint Spacing}: Filter waypoints to maintain a minimum spacing distance, reducing the total number of waypoints
	\end{enumerate}
	
	% ============================================================================
	\section{Implementation and Results}
	% ============================================================================
	
	\subsection{Project Structure}
	
	The implementation is organized in the following directory structure:
	
	\begin{lstlisting}[style=matlabstyle,caption={Project Structure}]
Perception_and_planning_lab/final_work/
|-- config/
|   |-- environment_map.m      % Map and obstacle configuration
|   |-- planner_parameters.m   % RRT* parameters
|-- planners/
|   |-- rrt_star.m             % Core RRT* algorithm
|   |-- generate_waypoints.m   % Waypoint generator wrapper
|-- utils/
|   |-- collision_check.m      % Collision detection
|   |-- path_smoothing.m       % Path optimization
|   |-- visualization_utils.m  % Plotting functions
|-- simulation/
|   |-- rrt_star_navigation_simulation.m  % Main simulation
|   |-- demo_rrt_star_planner.m           % RRT* demo
|-- tests/
|   |-- test_rrt_star_algorithm.m   % Unit tests
|   |-- test_full_navigation.m      % Integration tests
	\end{lstlisting}
	
	\subsection{Configuration Parameters}
	
	The main parameters that control the algorithm behavior are shown in Table~\ref{tab:params}.
	
	\begin{table}[htb]
		\centering
		\caption{RRT* Configuration Parameters}
		\begin{tabular}{lcc}
			\toprule
			\textbf{Parameter} & \textbf{Value} & \textbf{Description} \\
			\midrule
			max\_iterations & 4000 & Maximum tree expansion iterations \\
			step\_size & 0.8 m & Maximum extension distance \\
			goal\_bias & 10\% & Probability of sampling goal \\
			goal\_tolerance & 0.5 m & Distance to consider goal reached \\
			min\_waypoint\_spacing & 0.3 m & Minimum distance between waypoints \\
			\bottomrule
		\end{tabular}
		\label{tab:params}
	\end{table}
	
	\subsection{Available Scenarios}
	
	Multiple test scenarios are available:
	
	\begin{itemize}
		\item \texttt{empty}: Empty environment for basic testing
		\item \texttt{simple}: 4 obstacles (circles and rectangles)
		\item \texttt{maze}: Simple maze with walls
		\item \texttt{labyrinth}: Complete labyrinth with multiple routes
		\item \texttt{narrow}: Narrow passages
		\item \texttt{extreme}: 21 obstacles, corner-to-corner navigation
	\end{itemize}
	
	\subsection{Visualization Results}
	
	\begin{figure}[htpb]
		\centering
		\fbox{\parbox{0.8\textwidth}{\centering\vspace{3cm}\textit{[Insert RRT* tree expansion visualization here]}\vspace{3cm}}}
		\caption{RRT* tree expansion showing explored nodes and optimal path found.}
		\label{fig:rrt_tree}
	\end{figure}
	
	\begin{figure}[htpb]
		\centering
		\fbox{\parbox{0.8\textwidth}{\centering\vspace{3cm}\textit{[Insert path comparison: raw vs smoothed]}\vspace{3cm}}}
		\caption{Comparison between raw RRT* path and smoothed waypoints.}
		\label{fig:path_comparison}
	\end{figure}
	
	\begin{figure}[htpb]
		\centering
		\fbox{\parbox{0.8\textwidth}{\centering\vspace{3cm}\textit{[Insert robot navigation simulation screenshot]}\vspace{3cm}}}
		\caption{Real-time navigation simulation showing robot following RRT* generated waypoints.}
		\label{fig:navigation}
	\end{figure}
	
	\begin{figure}[htpb]
		\centering
		\fbox{\parbox{0.8\textwidth}{\centering\vspace{3cm}\textit{[Insert labyrinth scenario result]}\vspace{3cm}}}
		\caption{Path planning result in the labyrinth scenario with multiple possible routes.}
		\label{fig:labyrinth}
	\end{figure}
	
	% ============================================================================
	\section{How to Execute the Code}
	% ============================================================================
	
	\subsection{Prerequisites}
	
	\begin{itemize}
		\item MATLAB R2020a or later
		\item Fuzzy Logic Toolbox (for the robot controller)
		\item The complete \texttt{RoboticsLab\_DRM} repository
	\end{itemize}
	
	\subsection{Execution Steps}
	
	\subsubsection{Step 1: Navigate to Simulation Directory}
	
	\begin{lstlisting}[style=matlabstyle,caption={Navigate to simulation folder}]
% In MATLAB Command Window:
cd Perception_and_planning_lab/final_work/simulation
	\end{lstlisting}
	
	\subsubsection{Step 2: Run the Simulation}
	
	\begin{lstlisting}[style=matlabstyle,caption={Execute main simulation}]
% Run with default scenario (labyrinth):
rrt_star_navigation_simulation

% Or specify a scenario:
rrt_star_navigation_simulation('simple')
rrt_star_navigation_simulation('maze')
rrt_star_navigation_simulation('labyrinth')
rrt_star_navigation_simulation('extreme')
	\end{lstlisting}
	
	\subsubsection{Step 3: Observe the Results}
	
	The simulation will:
	\begin{enumerate}
		\item Generate the RRT* path with real-time tree visualization
		\item Display the smoothed waypoints
		\item Prompt to press any key to start navigation
		\item Show real-time robot navigation following the waypoints
		\item Display navigation statistics upon completion
	\end{enumerate}
	
	\subsubsection{Alternative: RRT* Demo Only}
	
	To visualize only the path planning algorithm without robot navigation:
	
	\begin{lstlisting}[style=matlabstyle,caption={Run RRT* demo}]
cd Perception_and_planning_lab/final_work/simulation
demo_rrt_star_planner
	\end{lstlisting}
	
	\subsubsection{Running Tests}
	
	\begin{lstlisting}[style=matlabstyle,caption={Execute tests}]
cd Perception_and_planning_lab/final_work/tests

% Unit tests for RRT* algorithm:
test_rrt_star_algorithm

% Integration test with robot controller:
test_full_navigation
	\end{lstlisting}
	
	% ============================================================================
	\section{Conclusion}
	% ============================================================================
	
	\justify
	The RRT* algorithm was successfully implemented and integrated with the existing skid-steer robot controller. The implementation demonstrates:
	
	\begin{itemize}
		\item Efficient path planning in complex environments with multiple obstacles
		\item Asymptotically optimal paths through the rewiring mechanism
		\item Real-time visualization of tree expansion and path optimization
		\item Seamless integration with the Fuzzy+PID+Lag Compensator controller
		\item Configurable parameters for different scenarios and requirements
	\end{itemize}
	
	The modular architecture allows easy extension to include dynamic obstacles, multi-robot coordination, or alternative planning algorithms in future work.
	
	% ============================================================================
	\newpage
	\appendix
	\section{Core Algorithm Code}
	% ============================================================================
	
	\begin{lstlisting}[style=matlabstyle,caption={RRT* Main Function (rrt\_star.m) - Excerpt}]
function [path, tree, stats] = rrt_star(start, goal, map, params)
% RRT* Path Planning Algorithm
% Finds asymptotically optimal path from start to goal

%% Initialize tree
tree.nodes = start';     % [x, y] for each node
tree.parent = 0;         % Parent index (0 = root)
tree.cost = 0;           % Cost from start

goal_found = false;
goal_idx = -1;

%% Main loop
for iter = 1:params.max_iterations
    % Sample random point (with goal bias)
    if rand() < params.goal_bias
        x_rand = goal;
    else
        x_rand = [map.x_min + rand()*(map.x_max - map.x_min);
                  map.y_min + rand()*(map.y_max - map.y_min)];
    end
    
    % Find nearest node
    [~, nearest_idx] = min(vecnorm(tree.nodes - x_rand', 2, 2));
    x_nearest = tree.nodes(nearest_idx, :)';
    
    % Steer toward random point
    direction = x_rand - x_nearest;
    dist = norm(direction);
    if dist > params.step_size
        direction = direction / dist * params.step_size;
    end
    x_new = x_nearest + direction;
    
    % Check collision
    if ~collision_check(x_nearest, x_new, map.obstacles, map.robot_radius)
        % Find neighbors within radius
        neighbor_radius = min(params.neighbor_radius_factor * ...
            sqrt(log(size(tree.nodes,1)+1) / (size(tree.nodes,1)+1)), ...
            params.step_size * 3);
        
        dists = vecnorm(tree.nodes - x_new', 2, 2);
        neighbor_idxs = find(dists < neighbor_radius);
        
        % Choose best parent (minimum cost)
        min_cost = tree.cost(nearest_idx) + norm(x_new - x_nearest);
        best_parent = nearest_idx;
        
        for n = neighbor_idxs'
            x_near = tree.nodes(n, :)';
            cost_through_n = tree.cost(n) + norm(x_new - x_near);
            if cost_through_n < min_cost
                if ~collision_check(x_near, x_new, map.obstacles, map.robot_radius)
                    min_cost = cost_through_n;
                    best_parent = n;
                end
            end
        end
        
        % Add new node
        tree.nodes = [tree.nodes; x_new'];
        tree.parent = [tree.parent; best_parent];
        tree.cost = [tree.cost; min_cost];
        new_idx = size(tree.nodes, 1);
        
        % Rewire neighbors
        for n = neighbor_idxs'
            x_near = tree.nodes(n, :)';
            cost_through_new = min_cost + norm(x_near - x_new);
            if cost_through_new < tree.cost(n)
                if ~collision_check(x_new, x_near, map.obstacles, map.robot_radius)
                    tree.parent(n) = new_idx;
                    tree.cost(n) = cost_through_new;
                end
            end
        end
        
        % Check if goal reached
        if norm(x_new - goal) < params.goal_tolerance
            goal_found = true;
            goal_idx = new_idx;
        end
    end
end

% Extract path by backtracking
if goal_found
    path = [];
    idx = goal_idx;
    while idx > 0
        path = [tree.nodes(idx, :); path];
        idx = tree.parent(idx);
    end
else
    path = [];
end
end
	\end{lstlisting}
	
\end{document}
